\SourceLeaf{42}

\begin{center}
\textbf{NACHLASS.}
\end{center}

\begin{center}
\textbf{ÜBER DIE INTEGRATION DER DIFFERENTIALFORMEL}\\
$(1+n\cos\varphi)^{\nu}. d\varphi.$
\end{center}

\noindent\hspace*{\parindent}\textls{Vorrede.} Der größte Teil der sogenannten Integralrechnung wird durch Entwicklung und Behandlung von Reihen bewältigt. Das einzige Verfahren, das man darin neben dem Probieren und dem Rückgriff auf bekannte Differentiationen findet, besteht darin, eine gegebene Differentialgröße in eine Reihe zu entwickeln, deren einzelne Glieder sich jeweils leicht integrieren lassen. Bei einer endlichen Reihe gelangt man so zu einem in jeder Hinsicht vollständigen Ausdruck. Eine unendliche Reihe stellt den gesuchten Zusammenhang der Unbekannten weniger vollkommen dar. Wenn sie jedoch so fortschreitet, dass eine Annäherung möglich ist, was sich in den meisten Fällen durch geeignete Behandlung erreichen lässt, genügt sie mit Recht für gewöhnliche Zwecke und gilt als gleichwertig mit einer Gleichung, sofern es zwischen den gegebenen Veränderlichen tatsächlich eine solche gibt. Wegen der meist umständlichen und scheinbar unregelmäßigen Glieder der Reihe ist es freilich eine schwierige Aufgabe, ihr Gesetz durch einen allgemeinen Ausdruck zu bestimmen und diesen Ausdruck durch einen Beweis zu sichern. Beides fordern die drohende Gefahr eines Fehlers, der sich in eine aus bloßer, plausibel gemachter Vermutung hervorgegangene Rechnung einschleichen kann, und die Würde der Wissenschaft. Die vorliegende kleine Abhandlung soll ein Beispiel einer solchen Untersuchung geben. Sie befasst sich mit der Integration der Formel $d\varphi(1+n\cos\varphi)^{\nu}$ durch Reihen, wobei $n$ und $\nu$ nach Belieben gewählte bestimmte Zahlen sind. Diese Formel ist von hervorragendem Nutzen in der physikalischen Astronomie und erschien dem berühmten \textsc{Euler} so wichtig, dass er ihr ohne Zögern ein ganzes Kapitel seiner Institutionum Calculi Integralis\HistoricalNote{*)}{*)}{Kap. VI, Bd. I.}%
